\documentclass[11pt]{article}
\usepackage[margin=1in]{geometry}
\usepackage{amsmath,amssymb,amsthm}
\usepackage{hyperref}
\usepackage[T1]{fontenc}
\usepackage[utf8]{inputenc}

\newtheorem{definition}{Definition}
\newtheorem{assumption}{Assumption}
\newtheorem{theorem}{Theorem}
\newtheorem{proposition}{Proposition}
\newtheorem{corollary}{Corollary}
\newtheorem{conjecture}{Conjecture}
\newtheorem{hypothesis}{Hypothesis}
\newtheorem{observation}{Observation}
\newtheorem{justification}{Justification}
\newtheorem{designprinciple}{Design Principle}
\newtheorem{example}{Example}

\title{OmegaSim/OmegaHive Strange-Attractor Tuning:\\
A Hyperseed-Dynamical Formalization}
\author{OmegaClaw / ZeroBot working notes for Ben Goertzel}
\date{2026-06-28}

\begin{document}
\maketitle

\begin{abstract}
This note formalizes a useful mathematical frame for constructing and tuning
OmegaSim/OmegaHive-style simulations so that they can explore complex bounded
nonlinear dynamics, including operational strange-attractor regimes.  The core
idea is to model hives as delayed, coupled, dissipative nonlinear systems; to
express tuning through dimensionless control ratios; to prove what follows from
well-posedness, boundedness, contraction, and delay-instability arguments; and to
translate the resulting theory into concrete simulation recommendations.  The
claims here are deliberately split into theorems, conjectures, operational
criteria, and implementation suggestions.
\end{abstract}

\section{Provenance and aim}

This note responds to Ben Goertzel's 2026-06-28 request for a rigorous
mathematical formalization of the ``useful frame'' proposed for tuning one-hive
and OmegaHive-style simulations toward complex strange-attractor dynamics.  The
requested steps are:
\begin{enumerate}
  \item formalize concepts bridging Hyperseed and nonlinear dynamics;
  \item state and prove general-purpose theorems applicable across agent hives;
  \item specialize these results to OmegaHive1/OmegaSim-style simulations;
  \item give theory-motivated instructions for the Codex/OmegaSim simulation
        agent.
\end{enumerate}

The purpose is not to prove that a given current simulation has a strange
attractor.  Rather, it is to identify mathematically meaningful design knobs and
verification gates so that such a claim can later become evidence-bearing.

\section{S1. Concepts bridging Hyperseed and nonlinear dynamics}

\subsection{Hyperseed event schema as a state observable}

\begin{definition}[Hyperseed process event]
A Hyperseed-relevant process event is a typed record
\[
  e = \langle \iota, \tau, A, C, R, P, \Theta, \nu \rangle
\]
where $\iota$ is identity, $\tau$ time, $A$ the active agents or sub-processes,
$C$ context, $R$ realized relation or result, $P$ provenance, $\Theta$ an
interpretive type assignment, and $\nu$ uncertainty/truth-value metadata.
\end{definition}

\begin{definition}[Hive observable]
A hive observable is a map
\[
  O : E_{\le t} \to X \subseteq \mathbb R^n
\]
from a finite or filtered history of Hyperseed process events to a numerical
state vector.  Components may include agent activation, resource load, attention,
coordination phase, memory pressure, trust, artifact quality, task backlog,
message volume, or semantic diversity.
\end{definition}

\begin{definition}[Scrutable reduction]
A reduction from Hyperseed events to dynamical state is scrutable when each
coordinate $x_i$ of $x=O(E)$ has: (i) a declared meaning, (ii) a measurement rule,
(iii) provenance to events or logs, and (iv) an uncertainty model.  Unscrutable
coordinates may be useful for machine learning, but should not carry ontological
interpretation without an additional explanation layer.
\end{definition}

\subsection{Delayed agent-hive dynamics}

Let $x(t)\in X\subseteq\mathbb R^n$ be a numerical hive state derived from the
observable map above.  Let the delayed history segment be
\[
  x_t(\theta)=x(t+\theta),\qquad \theta\in[-\tau,0],
\]
where $\tau\ge0$ is a maximum communication, reaction, or memory delay.

\begin{definition}[Continuous delayed hive system]
A continuous delayed hive system is a retarded functional differential equation
\[
  \dot x(t)=F(x_t,u(t),\eta(t);\theta),
\]
where $F:C([ -\tau,0],X)\times U\times E\to\mathbb R^n$, $u(t)$ is a controllable
intervention/tuning signal, $\eta(t)$ is bounded environmental or stochastic
forcing, and $\theta$ denotes fixed parameters.
\end{definition}

\begin{definition}[Discrete delayed hive simulation]
A discrete delayed hive simulation is a lifted map
\[
  x_{k+1}=G(x_k,x_{k-1},\ldots,x_{k-m},u_k,\eta_k;\theta),
\]
with lifted state
\[
  z_k=(x_k,x_{k-1},\ldots,x_{k-m})\in X^{m+1}.
\]
Equivalently,
\[
  z_{k+1}=\widehat G(z_k,u_k,\eta_k;\theta)
  =(G(x_k,\ldots,x_{k-m},u_k,\eta_k),x_k,\ldots,x_{k-m+1}).
\]
\end{definition}

\subsection{Dimensionless tuning coordinates}

The raw constants of a hive simulation are usually too numerous to sweep
intelligently.  We therefore define dimensionless ratios.

\begin{definition}[Dimensionless hive control ratios]
For a delayed nonlinear hive model, define:
\[
  \rho = g\,r(A),
  \qquad
  \delta = \frac{\tau}{T_{\rm relax}},
  \qquad
  \mu = \frac{T_{\rm memory}}{T_{\rm decision}},
\]
\[
  \kappa = \frac{\|A_{\rm cross}\|}{\|A_{\rm within}\|},
  \qquad
  \nu = \frac{\sigma_\eta}{\|F_0\|+\epsilon},
\]
where $g$ is a typical nonlinear response gain, $r(A)$ is the spectral radius of
an interaction matrix or linearized coupling operator, $T_{\rm relax}$ is local
relaxation time, $T_{\rm memory}$ memory decay time, $T_{\rm decision}$ update
time, $A_{\rm cross}$ and $A_{\rm within}$ are cross-hive and within-hive
couplings, and $\sigma_\eta$ is noise magnitude.
\end{definition}

Interpretation: $\rho$ measures feedback amplification, $\delta$ delay pressure,
$\mu$ persistence, $\kappa$ modular leakage, and $\nu$ noise-to-drift.  These are
not universal constants; they are organizing coordinates for phase diagrams.

\begin{definition}[Operational strange-attractor regime]
A parameter region $\mathcal P$ is an operational strange-attractor regime for a
simulation if, for a specified observable $O$ and lifted state $z$, the following
hold robustly over validation runs:
\begin{enumerate}
  \item trajectories are ultimately bounded and not merely saturating at hard
        numerical limits;
  \item the largest finite-time Lyapunov estimate is positive under controlled
        perturbations of initial histories;
  \item long-run statistics are stable under timestep/refinement and seed checks;
  \item recurrence plots, broadband spectra, or reconstructed dimensions indicate
        non-periodic structured recurrence rather than white noise.
\end{enumerate}
This is an empirical criterion.  It is weaker than a proof of a mathematical
strange attractor.
\end{definition}

\section{S2. General theorems}

\subsection{Well-posedness and boundedness}

\begin{theorem}[Local well-posedness for delayed hives]
Assume $F$ is locally Lipschitz in the history variable $x_t$, uniformly over
bounded $u$ and $\eta$.  Then for every continuous initial history
$\phi\in C([ -\tau,0],X)$ there exists a unique maximal solution.
\end{theorem}

\begin{proof}
This is the Picard--Lindelof theorem for retarded functional differential
equations.  The history space $C([ -\tau,0],X)$ is a Banach space under the sup
norm.  Local Lipschitz continuity gives a contraction on a sufficiently short
interval.  Repetition extends the solution until it exits every bounded subset of
the state space.
\end{proof}

\begin{theorem}[Ultimate boundedness from dissipativity]
Suppose there is a radially unbounded $C^1$ function
$V:\mathbb R^n\to\mathbb R_{\ge0}$ and constants $a,b>0$ such that along all
trajectories with bounded input magnitude $M$,
\[
  \dot V(x(t))\le -aV(x(t))+bM.
\]
Then
\[
  \limsup_{t\to\infty}V(x(t))\le \frac{bM}{a}.
\]
\end{theorem}

\begin{proof}
By Gronwall's inequality,
\[
  V(t)\le e^{-at}V(0)+\frac{bM}{a}(1-e^{-at}).
\]
Taking $t\to\infty$ gives the result.
\end{proof}

\begin{corollary}[Attractor-talk requires a bounded regime]
A simulation run that diverges, explodes numerically, or only bounces against
uninterpreted clipping bounds does not by itself justify talk of an attractor.
Dissipativity or empirically verified boundedness is a prerequisite diagnostic.
\end{corollary}

\begin{theorem}[Forward-invariant box for delayed maps]
For
\[
  x_{k+1}=G(x_k,\ldots,x_{k-m},u_k),
\]
assume there is a compact set $K\subseteq\mathbb R^n$ such that
\[
  G(K^{m+1},U)\subseteq K.
\]
Then $K^{m+1}$ is forward invariant under the lifted map $\widehat G$.
\end{theorem}

\begin{proof}
If $z_k=(x_k,\ldots,x_{k-m})\in K^{m+1}$, then each delayed component lies in $K$.
By assumption the new component $x_{k+1}=G(\cdots)$ lies in $K$.  The remaining
components of $z_{k+1}$ are shifted copies of previous components in $K$.
Therefore $z_{k+1}\in K^{m+1}$.
\end{proof}

\subsection{Contraction: when complex dynamics is impossible}

\begin{theorem}[Delayed-map contraction implies a unique attracting history]
Let $\widehat G:K^{m+1}\to K^{m+1}$ be the lifted delayed map on compact convex
$K^{m+1}$, and suppose it is a contraction with Lipschitz constant $L<1$ under
some norm.  Then there is a unique fixed lifted state $z^\ast$, and every
trajectory converges to $z^\ast$.
\end{theorem}

\begin{proof}
This is Banach's fixed-point theorem applied to $\widehat G$ on the complete
metric space $K^{m+1}$.  Iterates of a contraction converge to the unique fixed
point.
\end{proof}

\begin{corollary}[Low-gain hives cannot be chaotic]
If the effective lifted update of a hive simulation is contractive, then it cannot
have a strange attractor, quasiperiodic orbit, or sustained multi-cycle behavior.
Therefore one route toward complex dynamics is to increase gain, delay,
cross-coupling, or memory persistence until the lifted system leaves the
contraction regime while preserving boundedness.
\end{corollary}

\begin{proposition}[A practical Lipschitz bound]
Suppose
\[
  G(x_k,\ldots,x_{k-m})=S\left(B+\sum_{j=0}^m A_j x_{k-j}\right)
\]
where $S$ is componentwise Lipschitz with constant $s$ and the lifted norm is
chosen so that the shift has unit weight.  A sufficient condition for contraction
is that the induced lifted Lipschitz constant of the companion map is less than
one; in particular the response gain $s$ and delayed coupling norms $\|A_j\|$
provide an upper bound.  Thus the product ``gain times coupling'' is the first
order control axis.
\end{proposition}

\begin{proof}
The non-shift component changes by at most
$s\sum_j\|A_j\|\,\|\Delta x_{k-j}\|$ under the chosen induced norm, while the shift
components copy previous differences.  Weighted norms can penalize older delay
coordinates so that the whole companion map is contractive when the response and
coupling norms are sufficiently small.  This proves the sufficient, not necessary,
condition.
\end{proof}

\subsection{Delay-induced oscillatory instability}

\begin{theorem}[Scalar delay Hopf boundary]
Consider
\[
  \dot x(t)=-a x(t)+b x(t-\tau),\qquad a>0.
\]
The characteristic equation is
\[
  \lambda+a-b e^{-\lambda\tau}=0.
\]
A Hopf crossing candidate satisfies $\lambda=i\omega$, hence
\[
  a=b\cos(\omega\tau),\qquad
  \omega=-b\sin(\omega\tau).
\]
Therefore delayed feedback can destabilize an equilibrium through an oscillatory
root crossing even when the local undelayed term is dissipative.
\end{theorem}

\begin{proof}
Substitute $x(t)=e^{\lambda t}$ to obtain the characteristic equation.  At an
oscillatory stability boundary $\lambda=i\omega$.  Separating real and imaginary
parts gives the two displayed equations.
\end{proof}

\begin{corollary}[Why $\delta=\tau/T_{\rm relax}$ matters]
The ratio of communication/memory delay to local relaxation time is a genuine
bifurcation coordinate, not just an implementation detail.  Increasing delay can
convert damped consensus into oscillation and, in higher-dimensional nonlinear
systems, into period-doubling or more complex regimes.
\end{corollary}

\subsection{Existence of compact attractors}

\begin{theorem}[Compact global attractor under dissipative smoothing]
For a continuous delayed hive semiflow on $C([ -\tau,0],X)$, assume: (i) unique
forward solutions exist, (ii) solutions enter a bounded absorbing set, and (iii)
the semiflow is asymptotically compact.  Then there is a compact global attractor
$\mathcal A$ attracting all bounded sets.
\end{theorem}

\begin{proof}
This is the standard global-attractor theorem for continuous semiflows: a bounded
absorbing set plus asymptotic compactness implies the omega-limit of the absorbing
set is nonempty, compact, invariant, and attracts every bounded set.
\end{proof}

\begin{corollary}[Delay makes rigor harder]
Delayed systems are infinite-dimensional in history space.  Compactness and low
fractal dimension are not automatic.  OmegaSim should therefore treat finite
state-history simulations as approximations and validate attractor claims under
history-length and timestep variation.
\end{corollary}

\subsection{What remains conjectural}

\begin{conjecture}[Useful complex regime window]
For many OmegaHive-like systems there is a practical window between contraction
and blow-up/saturation, characterized roughly by $\rho\approx1$, non-negligible
$\delta$, moderate $\mu$, and controlled $\kappa$, in which bounded complex
recurrence is most likely.
\end{conjecture}

This is not a theorem because it depends on nonlinear form, topology, stochastic
forcing, and measurement.  It is a theory-guided search heuristic.

\begin{conjecture}[Hyperseed-scrutable complexity]
The most scientifically useful OmegaSim regimes will be those where the
low-dimensional observables show complex recurrence while Hyperseed event records
remain interpretable enough to explain which agent roles, memories, artifacts,
and feedback loops generate the recurrence.
\end{conjecture}

\section{S3. Specialization to OmegaHive1 / OmegaSim}

Since the exact OmegaHive1 simulation state is still evolving, define a generic
OmegaHive1 state vector as
\[
  x_k=(a_k,r_k,m_k,q_k,h_k,c_k,p_k)\in\mathbb R^n,
\]
where $a$ is agent activation/attention, $r$ resource or compute load, $m$ memory
persistence/pressure, $q$ artifact quality or progress, $h$ semantic diversity,
$c$ coordination/coherence, and $p$ task pressure.  The components may be vectors
indexed by agents, roles, artifacts, or sub-hives.

A useful discrete OmegaHive1 skeleton is
\[
  x_{k+1}=S\left(W_0x_k+W_1x_{k-d}+W_m M_k+B u_k+\xi_k\right)-D x_k,
\]
with memory update
\[
  M_{k+1}=(1-\alpha)M_k+\alpha H(x_k,\text{events}_k),
\]
where $S$ is saturating nonlinear response, $W_0$ present-time coupling, $W_1$
delayed coupling, $M_k$ memory state, $u_k$ interventions, $\xi_k$ controlled
noise, and $D$ local decay.

\begin{proposition}[OmegaHive1 boundedness by saturation and resource accounting]
If each component of $S$ maps into a compact interval $K_i$ and resource/memory
updates are either projected into compact budgets or dissipative in the sense of
the boundedness theorem, then the lifted OmegaHive1 history has a compact forward
invariant set.
\end{proposition}

\begin{proof}
The saturating components satisfy the invariant-box theorem.  Resource and memory
coordinates either remain in explicit compact budgets or satisfy dissipative
bounds.  The product of these compact coordinate sets is compact and invariant
under the lifted delayed map.
\end{proof}

\begin{proposition}[OmegaHive1 contraction baseline]
If the fitted or designed Lipschitz constant of the lifted OmegaHive1 update is
below one, the system will converge to a unique stable regime and should not be
expected to exhibit strange-attractor-like dynamics.
\end{proposition}

\begin{proof}
Immediate from the delayed-map contraction theorem.
\end{proof}

\begin{proposition}[OmegaHive1 delay/gain destabilization target]
Let $x^\ast$ be an equilibrium of a smooth deterministic OmegaHive1 skeleton.
Linearizing gives
\[
  \delta x_{k+1}=J_0\delta x_k+J_1\delta x_{k-d}+J_m\delta M_k.
\]
The stability of $x^\ast$ is determined by the eigenvalues of the lifted companion
Jacobian.  Parameter sweeps that increase response gain, delayed coupling, or
memory persistence move these eigenvalues toward or beyond the unit circle.
\end{proposition}

\begin{proof}
Linearization of the lifted map gives an ordinary finite-dimensional map on
history and memory coordinates.  Local discrete-time stability is equivalent to
all eigenvalues of the companion Jacobian lying inside the unit disk.  Increasing
the norms of coupling and persistence terms changes this Jacobian continuously;
instability occurs when eigenvalues cross the unit circle.
\end{proof}

\subsection{OmegaHive1 order parameters}

For diagnosis, define low-dimensional order parameters:
\[
  R_k=\left|\frac{1}{N}\sum_{j=1}^N e^{i\phi_{j,k}}\right|
  \quad\text{(coordination/synchrony)},
\]
\[
  H_k=-\sum_a p_k(a)\log p_k(a)
  \quad\text{(behavioral or semantic diversity)},
\]
\[
  C_k=x_k^\top L x_k
  \quad\text{(network tension across graph Laplacian $L$)},
\]
\[
  Q_k=\text{artifact-progress or quality metric},
  \qquad
  B_k=\text{resource/backlog pressure}.
\]
Analyze trajectories in $(R,H,C,Q,B)$ as well as in raw state space.

\begin{designprinciple}[Do not search for chaos in raw logs first]
Raw agent traces are too high-dimensional and semantically mixed.  First define
scrutable observables and order parameters, then test whether the induced
observable dynamics has bounded recurrence, sensitivity, and nonperiodicity.
\end{designprinciple}

\section{S4. Theory-motivated OmegaSim instructions}

\subsection{Implementation targets}

\begin{enumerate}
  \item Implement the lifted delayed state $z_k=(x_k,\ldots,x_{k-m},M_k)$
        explicitly.  Do not hide delays in ad hoc buffers without logging them.
  \item Define a compact state budget.  Use saturating nonlinearities,
        normalization, or explicit resource conservation so runs cannot be
        mistaken for attractors while diverging.
  \item Expose the control ratios $\rho,\delta,\mu,\kappa,\nu$ as primary sweep
        parameters.  Raw constants should be derived from these ratios where
        possible.
  \item Compute the fixed point or empirical baseline regime at low gain.  Use it
        as the contraction/control case.
  \item Linearize the deterministic skeleton around equilibria or recurrent
        regimes.  Track eigenvalues of the lifted companion Jacobian as parameters
        change.
  \item Run sweeps that cross from $\rho<1$ to $\rho>1$, from $\delta\approx0$ to
        $\delta=O(1)$, and from weak to moderate memory persistence $\mu$.
  \item Near transitions, run hysteresis sweeps upward and downward in gain/delay
        to detect multistability and crises.
  \item Estimate finite-time Lyapunov exponents with paired trajectories, periodic
        renormalization, multiple perturbation directions, and matched noise
        paths.
  \item Add null controls: no-delay, linearized nonlinearity, degree-preserving
        rewiring, delay-shuffled, phase-randomized surrogate, and frozen-noise
        paired comparisons.
  \item Classify each region as fixed point, limit cycle, quasiperiodic,
        candidate chaotic, transient chaotic, noise-driven irregular, divergent,
        or trivial saturation.
\end{enumerate}

\subsection{Specific suggested sweep grid}

A first OmegaSim phase diagram should sweep:
\[
  \rho\in[0.2,3.0],\qquad
  \delta\in\{0,0.1,0.3,1,3\},
\]
\[
  \mu\in\{0.5,1,3,10\},\qquad
  \kappa\in\{0,0.05,0.2,0.8\},
  \qquad
  \nu\in\{0,10^{-3},10^{-2},10^{-1}\}.
\]
The exact ranges should be rescaled after measuring the relaxation time and
state norms of the implemented simulator.

\subsection{Acceptance criteria for a candidate strange-attractor region}

A region should only be reported as a candidate complex/strange-attractor regime
if it passes all of these gates:
\begin{enumerate}
  \item boundedness gate: norms and resource variables remain in a nontrivial
        compact range;
  \item refinement gate: the qualitative regime survives smaller timestep or
        longer history resolution;
  \item perturbation gate: positive finite-time Lyapunov estimate under fixed
        noise path;
  \item surrogate gate: behavior differs from no-delay, linearized, and
        phase-randomized controls;
  \item semantic gate: Hyperseed event records identify plausible feedback loops
        generating the dynamics;
  \item task gate: the regime is related to useful simulation behavior, not just
        aesthetic irregularity.
\end{enumerate}

\section{Conclusion}

The rigorous core is modest but useful.  If the hive update is ill-posed,
unbounded, or contractive, complex attractor claims are unsupported or impossible.
If the system is well-posed, bounded, non-contractive, delayed, and nonlinear,
then gain, delay, memory persistence, modular leakage, and noise become principled
coordinates for exploring bifurcations and candidate strange-attractor regimes.

The practical OmegaSim instruction is therefore: build the simulator so these
coordinates are explicit; enforce boundedness; sweep through the contraction
boundary; test delay-induced instability; and only call a regime
strange-attractor-like after Lyapunov, recurrence, surrogate, refinement, and
semantic-provenance checks pass.

% 2026-07-01 OmegaSim Feedback Revision
% Addresses: cognitive appraisal, latent motivation, semantic/artifact fields,
% costly prediction, hysteresis, residual-state lobe discovery, linear controls,
% and experiment families A6--A8.

\section*{Revision Addendum (2026-07-01): Cognitive Appraisal, Residual Lobes, and New Experiment Families}

This addendum responds to the 2026-07-01 feedback on the original note.  It
preserves all prior material and adds eight substantive extensions.

\subsection*{R1. Logistic Response as Bounded Thresholded Cognitive Appraisal}

The original note used a generic saturating nonlinearity $S$.  We now interpret
the nonlinear response as a \emph{cognitive appraisal function}: each agent or
hive component evaluates incoming signals through a bounded, thresholded
transform---not as arbitrary chaos injection, but as a principled encoding of
attention allocation, urgency, and resource commitment.

\begin{definition}[Cognitive appraisal response]
The componentwise response is a logistic (sigmoid) function
\[
  \sigma_\alpha(x;\theta,\beta)=\frac{1}{1+\exp(-\alpha(x-\theta))},
\]
where $\alpha>0$ is the appraisal steepness (sensitivity to deviations from
threshold $\theta$) and $\beta$ is the saturation level (asymptotic maximum
response).  The effective response in the hive update replaces $S(\cdot)$ with
\[
  S_i(\cdot)\;\mapsto\; \beta_i\,\sigma_{\alpha_i}(w_i^\top x_k + b_i;\theta_i,\beta_i),
\]
giving each component $i$ its own threshold, steepness, and amplitude.
\end{definition}

\begin{proposition}[Appraisal threshold as bifurcation coordinate]
The appraisal steepness $\alpha$ acts as a gain parameter: low $\alpha$ yields
near-linear, contractive dynamics; high $\alpha$ creates a switching regime with
sharp transitions.  The threshold $\theta$ shifts the operating point.  Together
$(\alpha,\theta)$ provide a two-parameter family of bifurcation diagrams that
interpolates between smooth contraction and thresholded multi-state dynamics.
\end{proposition}

\begin{proof}
For $\alpha\to 0$, $\sigma_\alpha$ approaches a constant linear function with
slope $\beta/4$, recovering the low-gain contractive regime.  For $\alpha\to\infty$,
$\sigma_\alpha$ approaches a Heaviside step, introducing discontinuity-class
bifurcations.  Intermediate $\alpha$ produces smooth but steep transitions whose
Jacobian norm scales as $\alpha\beta/4$ at the threshold, directly controlling
whether the lifted map is contractive.
\end{proof}

\subsection*{R2. Latent Motivational State with Explicit Dynamics}

\begin{definition}[Latent motivational state]
Introduce a latent motivational vector $\mu_k\in\mathbb R^p$ with explicit
dynamics
\[
  \mu_{k+1}=(1-\lambda)\mu_k+\lambda\,\Phi(x_k,\mu_k)+\zeta_k,
\]
where $\lambda\in(0,1)$ is the motivation update rate, $\Phi$ is a feedback
function from current hive state and current motivation to updated motivation,
and $\zeta_k$ is endogenous drive noise.  The motivation $\mu_k$ modulates the
appraisal thresholds:
\[
  \theta_i(\mu_k)=\theta_i^{(0)}+\gamma_i^\top\mu_k,
\]
so that motivational state shifts what deviations are attended to.
\end{definition}

\begin{assumption}[Motivational timescale separation]
We assume $\lambda\ll 1/T_{\rm decision}$, i.e., motivational state evolves
slower than the decision update.  This yields a slow--fast system where $x_k$
is the fast variable and $\mu_k$ is the slow variable, enabling
singular-perturbation analysis.
\end{assumption}

\begin{proposition}[Slow motivation induces bifurcation sweeping]
As $\mu_k$ evolves slowly, the appraisal thresholds $\theta_i(\mu_k)$ drift,
causing the fast dynamics' bifurcation parameters to sweep.  This can produce
intermittent visits to multiple attractor lobes, hysteresis across bifurcation
boundaries, and crisis-induced switching---all driven by internal motivational
state rather than external forcing.
\end{proposition}

\subsection*{R3. Separate Semantic and Artifact Fields and Their Coupling}

\begin{definition}[Semantic field]
The semantic field $s_k\in\mathbb R^{d_s}$ captures the representational content
flowing through the hive: message semantics, task descriptions, shared memory
content, and interpretive type assignments.  Its dynamics are
\[
  s_{k+1}=\Psi_s(s_k,x_k,\mu_k)+\xi^s_k,
\]
where $\Psi_s$ may involve embedding-space operations, semantic similarity
weighting, and type-dependent routing.
\end{definition}

\begin{definition}[Artifact field]
The artifact field $a_k\in\mathbb R^{d_a}$ tracks produced artifacts: code,
documents, test results, structured outputs.  Its dynamics are
\[
  a_{k+1}=\Psi_a(a_k,s_k,x_k,\mu_k)+\xi^a_k,
\]
where $\Psi_a$ includes quality evaluation, version tracking, and resource cost
accounting.
\end{definition}

\begin{definition}[Field coupling]
The coupled system replaces the monolithic state $x_k$ with
\[
  \tilde x_k=(x_k^{\rm op},s_k,a_k,\mu_k),
\]
where $x_k^{\rm op}$ are operational/coordination variables.  The semantic field
influences operational dynamics through appraisal thresholds $\theta_i(s_k)$,
and the artifact field feeds back through resource load $r_k=g(a_k)$ and
quality metrics $q_k=Q(a_k)$.  Bidirectional coupling strength is parameterized
by $\chi_{sa}$ (semantic$\to$artifact) and $\chi_{as}$ (artifact$\to$semantic).
\end{definition}

\begin{designprinciple}[Separate fields enable distinct diagnostics]
Treating $s_k$ and $a_k$ as distinct dynamical fields allows separate Lyapunov
analysis, attractor characterization, and surrogate controls for each.  Semantic
complexity without artifact complexity, or vice versa, is diagnostically distinct
from joint complexity.
\end{designprinciple}

\subsection*{R4. Costly Prediction and Resource Accounting}

\begin{definition}[Prediction cost]
Each predictive step---whether by an agent predicting outcomes or by the hive
forecasting its own state---incurs a resource cost
\[
  c_k^{\rm pred}=c_{\rm fixed}+c_{\rm var}\|p_k\|^2,
\]
where $p_k$ is the prediction depth or model complexity and $c_{\rm var}>0$.
Total resource consumption must satisfy
\[
  \sum_{i} c_{k,i}^{\rm pred}+c_{k,i}^{\rm act}\le R_{\max},
\]
where $R_{\max}$ is the hive-wide resource budget per timestep.
\end{definition}

\begin{definition}[Resource-constrained dynamics]
The hive update with resource accounting becomes
\[
  x_{k+1}=S\bigl(W_0x_k+W_1x_{k-d}+W_mM_k+B u_k+\xi_k;\,R_{\max}-c_k^{\rm pred}\bigr),
\]
where the remaining budget $R_{\max}-c_k^{\rm pred}$ gates the effective
amplitude of activation.  When prediction costs are high, effective gain drops,
pushing the system toward the contractive regime; when prediction is cheap,
full nonlinear dynamics operate.
\end{definition}

\begin{proposition}[Cost-induced self-regulation]
If prediction cost scales with the complexity of the dynamics (e.g.,
$c_k^{\rm pred}\propto$ local finite-time Lyapunov exponent), then highly
chaotic regimes automatically incur higher costs, reducing effective gain and
creating a self-regulating feedback loop.  This provides an intrinsic mechanism
for bounded complexity without external clipping.
\end{proposition}

\subsection*{R5. Hysteresis and Adaptive Thresholds}

\begin{definition}[Adaptive appraisal thresholds]
Appraisal thresholds evolve according to
\[
  \theta_{i,k+1}=\theta_{i,k}+\eta_i\bigl(\sigma_{\alpha_i}(w_i^\top x_k;\theta_{i,k})-p_i^{\ast}\bigr),
\]
where $p_i^{\ast}\in(0,1)$ is a target activation rate and $\eta_i>0$ is the
adaptation rate.  This drives each component toward a target operating point,
implementing homeostatic regulation.
\end{definition}

\begin{proposition}[Adaptive thresholds create hysteresis]
When the adaptation rate $\eta_i$ is small but nonzero, the threshold
$\theta_{i,k}$ tracks the recent activation level.  On upward parameter sweeps
(increasing gain $\alpha$), the threshold rises to suppress over-activation; on
downward sweeps, it lowers slowly.  This asymmetric adaptation produces
hysteresis: the system's state at gain $\alpha$ on the upward sweep differs from
its state at the same $\alpha$ on the downward sweep.  The hysteresis loop width
scales as $O(\eta_i / \lambda_{\rm bif})$ where $\lambda_{\rm bif}$ is the
bifurcation parameter sweep rate.
\end{proposition}

\begin{proof}
Let $\bar\theta(\alpha)$ be the quasi-steady threshold satisfying
$\sigma_{\alpha}(w^\top x^\ast(\alpha,\bar\theta);\bar\theta)=p^\ast$.
Under slow sweep $\alpha_k=\alpha_0+\epsilon k$, the actual threshold lags:
$\theta_{k}\approx\bar\theta(\alpha_k)-\frac{\eta_i}{\epsilon}\bar\theta'(\alpha_k)$.
The correction term is direction-dependent (it changes sign with sweep
direction), producing a hysteresis gap proportional to $\eta_i/\epsilon$.
\end{proof}

\subsection*{R6. Residual-State Lobe Discovery}

\begin{definition}[Residual state]
Define the residual state as the deviation of the full system state from the
best-fit low-dimensional attractor model:
\[
  \rho_k = \tilde x_k - \Pi_{\mathcal A}(\tilde x_k),
\]
where $\tilde x_k=(x_k^{\rm op},s_k,a_k,\mu_k)$ is the full coupled state,
$\Pi_{\mathcal A}$ is the projection onto the fitted attractor model
$\mathcal A$ (e.g., a principal manifold, delay-coordinate embedding, or
low-order DMD approximation), and $\rho_k$ is the residual.  The residual
captures structure not explained by the low-dimensional attractor model.
\end{definition}

\begin{definition}[Residual lobe]
A residual lobe is a persistent cluster in residual-state space
$\{\rho_k\}$ identified by a clustering algorithm (e.g., Gaussian mixture,
HDBSCAN, or $k$-medoids in delay-coordinate space) with the following
properties:
\begin{enumerate}
  \item \emph{Persistence}: the cluster survives across multiple validation
        windows of length $\ge 10\,T_{\rm relax}$.
  \item \emph{Separation}: the cluster's centroid is at least $2\sigma$ from
        other cluster centroids in residual norm.
  \item \emph{Reproducibility}: the cluster appears in $\ge 80\%$ of
        independent runs with different initial histories under the same
        parameter setting.
  \item \emph{Semantic distinguishability}: Hyperseed event records for
        trajectories within the lobe differ systematically (in agent roles,
        artifact types, or semantic content) from records for other lobes.
\end{enumerate}
\end{definition}

\begin{designprinciple}[Lobe discovery as a validation tool]
Residual lobe discovery provides a concrete method for detecting whether the
full system has multi-lobe attractor structure that is invisible to the
low-dimensional observable.  If lobes are semantically distinguishable, they
identify distinct cognitive regimes (e.g., exploration vs.\ exploitation,
problem-attacking vs.\ reflective) that the hive transitions between.
\end{designprinciple}

\subsection*{R7. Amplitude-Matched Linear Controls}

\begin{definition}[Amplitude-matched linear control]
Given a nonlinear hive system $x_{k+1}=G(x_k,\ldots,x_{k-m})$ with operational
state norm $\|x_k\|\le X_{\max}$ in the regime of interest, the
amplitude-matched linear control is
\[
  x_{k+1}^{\rm lin}=A_0 x_k + A_1 x_{k-d} + B u_k + \bar c,
\]
where $A_0,A_1$ are obtained by linearizing $G$ around the regime mean
$\bar x = \frac{1}{T}\sum_{k} x_k$ and $\bar c$ is chosen so that
$\|x_k^{\rm lin}\|\approx\|x_k\|$ in the matched regime.  The linear control
uses the same noise path $\xi_k$ (common random numbers) and the same delay
structure as the nonlinear system.
\end{definition}

\begin{proposition}[Linear control as a surrogate baseline]
The amplitude-matched linear control isolates the contribution of nonlinearity
from delay and noise.  If the linear control exhibits comparable spectral
breadth, Lyapunov exponent, or recurrence statistics, the nonlinear system's
complexity is attributable to delay/noise, not to nonlinear structure.  Only
when the nonlinear system significantly exceeds the linear control on these
metrics can nonlinear strange-attractor claims be supported.
\end{proposition}

\begin{proof}
The linear control shares delay structure, noise realization, and amplitude
with the nonlinear system; it differs only in the replacement of $G$ by its
linearization.  Therefore any excess complexity in the nonlinear system is
attributable to the nonlinear terms.  This is a standard matched-control
argument from the design-of-experiments tradition.
\end{proof}

\subsection*{R8. Experiment Families A6, A7, and A8}

\subsubsection*{A6: Cognitive Appraisal Bifurcation Scan}

\textbf{Hypothesis:} Varying appraisal steepness $\alpha$ and threshold
$\theta$ across the hive produces a smooth bifurcation from contractive
fixed-point dynamics through damped oscillation to multi-lobe recurrent
dynamics, mediated by the cognitive appraisal function rather than by external
noise injection.

\textbf{Sweep variables:}
\begin{itemize}
  \item Appraisal steepness $\alpha\in\{0.5,1,2,5,10,20\}$
  \item Threshold $\theta\in\{-2,-1,0,1,2\}$ (relative to baseline input
        distribution)
  \item Motivational coupling $\gamma\in\{0,0.1,0.5,1.0\}$
\end{itemize}

\textbf{Controls:}
\begin{itemize}
  \item Amplitude-matched linear control (R7) at each $(\alpha,\theta)$ point
  \item No-motivational-dynamics control ($\gamma=0$, frozen $\mu_k$)
  \item Phase-randomized surrogate
\end{itemize}

\textbf{Acceptance criteria:}
\begin{enumerate}
  \item Bifurcation diagram in $(\alpha,\theta)$ shows a qualitative
        transition from fixed-point to oscillatory to multi-lobe dynamics.
  \item The transition region has positive finite-time Lyapunov exponent
        while the linear control does not.
  \item Motivational coupling ($\gamma>0$) widens the multi-lobe region
        compared to $\gamma=0$.
  \item Hysteresis is observable on upward vs.\ downward $\alpha$ sweeps with
        adaptive thresholds.
\end{enumerate}

\subsubsection*{A7: Residual Lobe Discovery and Validation}

\textbf{Hypothesis:} The full coupled system $(x^{\rm op},s,a,\mu)$ has
multi-lobe attractor structure that is partially invisible to the low-dimensional
observable $O(E_{\le t})$.  Residual lobes are semantically distinguishable.

\textbf{Sweep variables:}
\begin{itemize}
  \item Attractor model dimension $d_{\mathcal A}\in\{2,3,5,8\}$
  \item Clustering algorithm: HDBSCAN, GMM, $k$-medoids
  \item Delay embedding dimension for residual space: $m_{\rm emb}\in\{3,5,10\}$
  \item Validation window length: $\{5,10,20\}\cdot T_{\rm relax}$
  \item Number of independent runs: $\ge 20$ per parameter setting
\end{itemize}

\textbf{Controls:}
\begin{itemize}
  \item Random-label surrogate: cluster assignments randomly permuted across
        time to test whether semantic distinguishability exceeds chance.
  \item Single-field controls: residual computed using only $x^{\rm op}$, only
        $s$, only $a$, or only $\mu$, to test which field's residual carries
        lobe structure.
  \item Null-model residual: residual computed against a trivial (zero-order)
        attractor model, i.e., the mean state, to verify that the clustering
        captures genuine attractor structure and not just noise.
\end{itemize}

\textbf{Acceptance criteria:}
\begin{enumerate}
  \item At least one residual lobe satisfies all four properties from the
        residual lobe definition (R6): persistence, separation, reproducibility,
        and semantic distinguishability.
  \item The random-label surrogate produces $<5\%$ false-positive lobe rate
        at the same clustering thresholds.
  \item Single-field controls identify which field(s) contribute most to lobe
        structure, enabling targeted interpretation.
  \item Lobe membership predicts subsequent behavioral regime with above-chance
        accuracy ($p<0.01$ by permutation test).
\end{enumerate}

\subsubsection*{A8: Costly Prediction and Self-Regulation Experiment}

\textbf{Hypothesis:} When prediction cost scales with dynamical complexity
(high local Lyapunov exponent $\Rightarrow$ higher $c_k^{\rm pred}$), the hive
system self-regulates into a bounded-complexity regime: chaotic episodes are
automatically damped by resource depletion, and the system settles into an
intermittent regime alternating between high-complexity bursts and
resource-recovery quiescence.

\textbf{Sweep variables:}
\begin{itemize}
  \item Prediction cost coefficient $c_{\rm var}\in\{0,0.01,0.1,0.5,1.0\}$
  \item Resource budget $R_{\max}\in\{0.5,1,2,5\}$ (relative to median
        unconstrained activation cost)
  \item Cost complexity coupling: $c_k^{\rm pred}\propto |\lambda_{\rm FTLE}|^p$
        with exponent $p\in\{0,1,2\}$ (no coupling, linear, quadratic)
  \item Prediction depth $p_k\in\{1,2,5,10\}$ steps ahead
\end{itemize}

\textbf{Controls:}
\begin{itemize}
  \item Zero-cost control ($c_{\rm var}=0$): standard dynamics without
        resource accounting, to establish the unconstrained complexity baseline.
  \item Fixed-cost control: $c_k^{\rm pred}=c_{\rm fixed}$ (constant regardless
        of complexity), to distinguish resource scarcity effects from
        complexity-coupled self-regulation.
  \item Amplitude-matched linear control (R7) with the same resource budget,
        to separate nonlinear self-regulation from linear resource gating.
\end{itemize}

\textbf{Acceptance criteria:}
\begin{enumerate}
  \item With cost-complexity coupling ($p>0$), the system exhibits
        intermittent complexity: episodes of positive finite-time Lyapunov
        exponent alternated with quiescent recovery periods.
  \item The zero-cost control shows sustained chaos (no self-damping),
        confirming that the intermittent pattern is caused by cost coupling,
        not by other parameters.
  \item The fixed-cost control shows uniform suppression (no intermittency),
        confirming that the intermittency requires complexity-dependent cost,
        not just resource scarcity.
  \item Long-run complexity (time-averaged Lyapunov exponent) decreases
        monotonically with $c_{\rm var}$ for fixed $R_{\max}$, demonstrating
        a tunable self-regulation knob.
  \item Hysteresis in the complexity--cost relationship is observable when
        $c_{\rm var}$ is swept up and down, due to adaptive thresholds (R5).
\end{enumerate}

\subsection*{Summary of revisions}

This addendum adds: (R1) cognitive appraisal interpretation of the logistic
response, (R2) latent motivational state with slow--fast dynamics and
threshold modulation, (R3) separate semantic and artifact fields with
bidirectional coupling, (R4) costly prediction with resource accounting and
self-regulating complexity, (R5) adaptive thresholds producing hysteresis,
(R6) residual-state lobe discovery with clustering and validation, (R7)
amplitude-matched linear controls isolating nonlinear contributions, and (R8)
three new experiment families (A6, A7, A8) with explicit hypotheses, sweep
variables, controls, and acceptance criteria.  All original material is
preserved; this content is appended via \texttt{\\input} before
\texttt{\\end{document}}.

\end{document}
